\startSection
\selectlanguage{portuguese}
	\sectionTitle{Hino Nacional Brasileiro}
	
	\begin{question}{Considerando o processo de oficialização do Hino Nacional Brasileiro no final do século XIX, qual compositor foi responsável pela criação da melodia que se tornou oficialmente reconhecida em 1890?}
		\choice{Carlos Gomes}
		\choice{Alberto Nepomuceno}
		\choice[correct]{Francisco Manuel da Silva}
		\choice{Heitor Villa-Lobos}
		\choice{Antônio Carlos Jobim}
	\end{question}

	
	\begin{question}{O Hino Nacional Brasileiro passou por diferentes fases. Em que data ele foi executado pela primeira vez, ainda sob o título de \enquote{Hino ao 7 de abril}, em homenagem à abdicação de Dom Pedro I?}
		\choice{7 de setembro de 1822}
		\choice{15 de novembro de 1889}
		\choice{21 de abril de 1831}
		\choice{25 de março de 1824}
		\choice[correct]{13 de abril de 1831}
	\end{question}

	
	\begin{question}{Durante as comemorações do centenário da Independência do Brasil, foi composta a letra que acompanha o Hino Nacional Brasileiro até hoje. Quem foi o autor que se tornou parte oficial do hino?}
		\choice{Olavo Bilac}
		\choice{Gonçalves Dias}
		\choice{Castro Alves}
		\choice[correct]{Joaquim Osório Duque-Estrada}
		\choice{Manuel Bandeira}
	\end{question}

	
	\begin{question}{De acordo com a Lei n° 5.700/1971, que regulamenta os símbolos nacionais, em qual das situações abaixo é exigida a execução apenas da introdução e do coro do Hino Nacional Brasileiro, sem a repetição da segunda parte?}
		\choice[correct]{Cerimônias religiosas de sentido patriótico}
		\choice{Abertura de sessões cívicas}
		\choice{Hasteamento da Bandeira Nacional}
		\choice{Cerimônias oficiais do Governo Federal}
		\choice{Programas de rádio e televisão às 19 horas}
	\end{question}

	
	\begin{question}{O Hino Nacional Brasileiro contém diversas imagens poéticas que exaltam o país. Qual dos versos abaixo faz uma referência direta à constelação do Cruzeiro do Sul, símbolo presente na bandeira nacional?}
		\choice{\enquote{Gigante pela própria natureza}}
		\choice{\enquote{Fulguras, ó Brasil, florão da América}}
		\choice[correct]{\enquote{A imagem do Cruzeiro resplandece}}
		\choice{\enquote{Em teu formoso céu, risonho e límpido}}
		\choice{\enquote{Do que a terra mais garrida}}
	\end{question}
\renderSection
